% Chapter 2 — Tensors and Differential Forms (reviewed)
\chapter{Tensors and Differential Forms}

There are a number of vector spaces and algebras naturally associated with the tangent space  $ M_{m} $. Suitably smooth assignments of elements of these spaces to the points in M yield tensor fields and differential forms of various types. We shall first develop some of the pertinent facts from multilinear algebra, and then beginning with 2.14 we shall apply these concepts to manifolds.

\section{Tensor and Exterior Algebras}


Throughout 2.1–2.13, $V$, $W$, and $U$ will denote finite dimensional real vector spaces. As usual, $V^{*}$ will denote the dual space of $V$ consisting of all real-valued linear functions on $V$.

\warnernumber{1}
\begin{definition}\label{thm:2.1}
Let $F(V,W)$ be the free vector space over $\mathbb{R}$ whose generators are the points of $V \times W$. Thus $F(V,W)$ consists of all finite linear combinations of pairs $(v,w)$ with $v \in V$ and $w \in W$. Let $R(V,W)$ be the subspace of $F(V,W)$ generated by the set of all elements of $F(V,W)$ of the following forms:

\begin{equation}
\tag{1}\label{eq:2.1:1}
\begin{array}{l}{(v_{1}+v_{2},w)-(v_{1},w)-(v_{2},w)}\\ {(v,w_{1}+w_{2})-(v,w_{1})-(v,w_{2})}\\ {(a v,w)-a(v,w)}\\ {(v,a w)-a(v,w)}\\ \end{array}\left(\begin{array}{c}{a\in\mathbb{R}}\\ {v,v_{1},v_{2}\in V}\\ {w,w_{1},w_{2}\in W}\\ \end{array}\right).
\end{equation}

The quotient space $F(V,W)/R(V,W)$ is called the tensor product of $V$ and $W$ and is denoted by $V\otimes W$. The coset of $V\otimes W$ containing the element $(v,w)$ of $F(V,W)$ is denoted by $v\otimes w$. It follows from (1) that we have the following identities in $V\otimes W$:

\[
\begin{aligned}{(v_{1}+v_{2})\otimes w}&{{}=v_{1}\otimes w+v_{2}\otimes w}\\ {v\otimes(w_{1}+w_{2})}&{{}=v\otimes w_{1}+v\otimes w_{2}}\\ {a(v\otimes w)}&{{}=a v\otimes w=v\otimes a w.}\\ \end{aligned}
\]

\end{definition}

\warnernumber{2}
\discuss{2.2}{}\label{thm:2.2}
The following properties of the tensor product are easily established, and are left to the reader as exercises.

(a) Universal Mapping Property. Let $\varphi$ denote the bilinear map $(v,w)\mapsto$ $v\otimes w$ of $V\times W$ into $V\otimes W$. Then whenever $U$ is a vector space and $l:V\times W\to U$ is a bilinear map, there exists a unique linear map $\tilde{l}\colon V\otimes W\to U$ such that the following diagram commutes:


The pair consisting of $V\otimes W$ and $\varphi$ is said to solve the universal mapping problem for bilinear maps with domain $V\times W$. Moreover, $V\otimes W$ and $\varphi$ are unique with this property in the sense that if $X$ is a vector space and $\tilde{\varphi}:V\times W\to X$ a bilinear map with the above universal mapping property, then there exists an isomorphism $\alpha:V\otimes W\to X$ such that $\alpha\circ\varphi=\tilde{\varphi}$.

\begin{center}
\redrawn{figure-22}
\end{center}

(b)  $ V \otimes W $ is canonically isomorphic with  $ W \otimes V $.

(c)  $ V \otimes (W \otimes U) $ is canonically isomorphic with  $ (V \otimes W) \otimes U $.

(d) By property (a), the bilinear map of $V^{*}\times W$ into the vector space $\mathrm{Hom}(V,W)$ of linear transformations from $V$ to $W$ defined by $(f,w)(v)=f(v)\cdot w$ for $f\in V^{*}$, $v\in V$, and $w\in W$ determines uniquely a linear map $\alpha\colon V^{*}\otimes W\to\mathrm{Hom}(V,W)$. $\alpha$ is an isomorphism. As a consequence,


\[
\dim V\otimes W=(\dim V)(\dim W).
\]

(e) Let $\{e_i: i = 1, \ldots, c\}$ and $\{f_j: j = 1, \ldots, d\}$ be bases for $V$ and $W$ respectively. Then $\{e_i \otimes f_j: i = 1, \ldots, c$ and $j = 1, \ldots, d\}$ is a basis of $V \otimes W$.


\warnernumber{3}
\begin{definition}\label{thm:2.3}
The tensor space $V_{r,s}$ of type $(r,s)$ associated with $V$ is the vector space


\[
\underbrace{V\otimes\cdots\otimes V}_{r\text{ copies}}\otimes\underbrace{V^{*}\otimes\cdots\otimes V^{*}}_{s\text{ copies}}.
\]

The direct sum


\[
T(V)=\sum V_{r,s}\qquad(r,s\geq0),
\]

where  $ V_{0,0} = \mathbb{R} $, is called the tensor algebra of  $ V $. Elements of  $ T(V) $ are finite linear combinations over  $ \mathbb{R} $ of elements of the various  $ V_{r,s} $ and are called tensors.  $ T(V) $ is a non-commutative, associative, graded algebra under  $ \otimes $ multiplication: if  $ u = u_1 \otimes \cdots \otimes u_{r_1} \otimes u_1^* \otimes \cdots \otimes u_{s_1}^* $ belongs to  $ V_{r_1,s_1} $ and  $ v = v_1 \otimes \cdots \otimes v_{r_2} \otimes v_1^* \otimes \cdots \otimes v_{s_2}^* $ belongs to $V_{r_2,s_2}$, \index{Tensor algebra}\index{Tensors}




then their product  $ u \otimes v $ is defined by

\[
u\otimes v=u_{1}\otimes\cdots\otimes u_{r_{1}}\otimes v_{1}\otimes\cdots\otimes v_{r_{2}}\otimes u_{1}^{*}\otimes\cdots\otimes u_{s_{1}}^{*}\otimes v_{1}^{*}\otimes\cdots\otimes v_{s_{2}}^{*}
\]

and belongs to $V_{r_{1}+r_{2},s_{1}+s_{2}}$. Tensors in a particular tensor space $V_{r,s}$ are called homogeneous of degree $(r,s)$. A homogeneous tensor (of degree $(r,s)$, say) is called decomposable if it can be written in the form

\[
v_{1}\otimes\cdots\otimes v_{r}\otimes v_{1}^{*}\otimes\cdots\otimes v_{s}^{*}
\]

where  $ v_i \in V $  $ (i = 1, \ldots, r) $ and  $ v_j^* \in V^* $  $ (j = 1, \ldots, s) $.
\end{definition}

\warnernumber{4}
\begin{definition}\label{thm:2.4}
We let $C(V)$ denote the subalgebra $\sum_{k=0}^{\infty} V_{k,0}$ of $T(V)$. Let $I(V)$ be the two-sided ideal in $C(V)$ generated by the set of elements of the form $v\otimes v$ for $v\in V$, and set

\begin{equation}
\tag{1}\label{eq:2.4:1}
I_{k}(V)=I(V)\cap V_{k,0}.
\end{equation}

It follows that

\begin{equation}
\tag{2}\label{eq:2.4:2}
\begin{array}{r}{I(V)=\displaystyle\sum_{k=0}^{\infty}I_{k}(V),}\end{array}
\end{equation}

and is a graded ideal in $C(V)$. The exterior algebra $\Lambda(V)$ of $V$ is the graded algebra $C(V)/I(V)$. If we set \index{Exterior algebra}

\begin{equation}
\tag{3}\label{eq:2.4:3}
\Lambda_{k}(V)=V_{k,0}/I_{k}(V)\qquad(k\geq2),\qquad\Lambda_{0}(V)=\mathbb{R},\qquad\Lambda_{1}(V)=V,
\end{equation}

then

\begin{equation}
\tag{4}\label{eq:2.4:4}
\Lambda(V)=\sum_{k=0}^{\infty}\Lambda_{k}(V).
\end{equation}

We shall denote multiplication in the algebra $\Lambda(V)$ by $\wedge$. This is called the wedge or exterior product. In particular, the residue class containing $v_1\otimes\cdots\otimes v_k$ is $v_1\wedge\cdots\wedge v_k$. \index{Exterior (wedge) product}

\end{definition}

\warnernumber{5}
\begin{definition}\label{thm:2.5}
A multilinear map

\[
h\colon\underbrace{V\times\cdots\times V}_{r\text{ copies}}\to W
\]

is called alternating if

\[
h(v_{\pi(1)},\ldots,v_{\pi(r)})=(\operatorname{sgn}\pi)h(v_{1},\ldots,v_{r})\qquad(v_{1},\ldots,v_{r}\in V),
\]

for all permutations $\pi$ in the permutation group $S_{r}$ on $r$ letters. Sgn $\pi$ is the sign of the permutation $\pi$ (+1 if $\pi$ is even, -1 if $\pi$ is odd). The vector space of all alternating multilinear functions

\[
\underbrace{V\times\cdots\times V}_{r\text{ copies}}\to\mathbb{R}
\]

will be denoted by  $ A_r(V) $, and for convenience we set  $ A_0(V) = \mathbb{R} $.

\end{definition}

\warnernumber{6}
\discuss{2.6}{}\label{thm:2.6}
The following properties of the exterior algebra are left to the reader as exercises:

(a)  If $u \in \Lambda_k(V) $ and  $ v \in \Lambda_l(V) $, then  $ u \land v \in \Lambda_{k+l}(V) $ and  $ u \land v = (-1)^{kl} v \land u $.

(b) If  $ e_{1}, \ldots, e_{d} $ is a basis of V, then


\begin{equation}
\tag{1}\label{eq:2.6:1}
\{e_{\Phi}\}
\end{equation}

is a basis of $\Lambda(V)$, where $\Phi$ runs over all subsets of $\{1, \ldots, d\}$, including the empty set; where $e_{\Phi} = e_{i_1} \wedge \cdots \wedge e_{i_r}$ with $i_1 < \cdots < i_r$ when $\Phi$ is the subset $\{i_1, \ldots, i_r\}$ of $\{1, \ldots, d\}$; and where $e_{\Phi} = 1$ when $\Phi = \varnothing$. In particular,

\[
\Lambda_{d}(V)\cong\mathbb{R},
\]

\begin{equation}
\tag{2}\label{eq:2.6:2}
\Lambda_{d+j}(V)=\{0\}\qquad(j>0).
\end{equation}

Moreover, it follows that

\[
\dim\Lambda(V)=2^{d},
\]

\begin{equation}
\tag{3}\label{eq:2.6:3}
\dim\Lambda_{k}(V)=\binom{d}{k}=\frac{d!}{k!(d-k)!}\qquad(0\leq k\leq d).
\end{equation}

(Hint: Observe that the elements $\{e_\Phi\}$ span $\Lambda(V)$. To prove that they also are linearly independent, first prove that $e_1 \wedge \cdots \wedge e_d$ is not zero in $\Lambda_d(V)$. For this, one must show that $e_1 \otimes \cdots \otimes e_d$ does not belong to $I(V)$. Express an arbitrary element of $I(V)$ in terms of the basis vectors $e_1, \ldots, e_d$, and show that it could not equal $e_1 \otimes \cdots \otimes e_d$. Then for linear independence of the entire set $\{e_\Phi\}$, multiply the equation $\sum a_e e_\Phi = 0$ by suitable products of the $e_i$ to land in $\Lambda_d(V)$, and conclude that the various $a_\Phi$ are all zero.)

(c) Universal Mapping Property. Let $\varphi$ denote the mapping $(v_1, \ldots, v_k) \mapsto v_1 \wedge \cdots \wedge v_k$ of $V \times \cdots \times V$ ($k$ copies) into $\Lambda_k(V)$. Then $\varphi$ is an alternating multilinear map. Now to each alternating multilinear map $h$ of $V \times \cdots \times V$ ($k$ copies) into a vector space $W$, there corresponds uniquely a linear map $\tilde{h}$: $\Lambda_k(V) \to W$ such that $\tilde{h} \circ \varphi = h$. \index{Universal mapping property of exterior algebra}


\begin{equation}
\tag{4}\label{eq:2.6:4}\redrawn{figure-23}
\end{equation}

The pair consisting of $\Lambda_{k}(V)$ and $\varphi$ is said to solve the universal mapping problem for alternating multilinear maps with domain $V\times\cdots\times V$; and this is the unique solution in the sense that if $X$ is a vector space and $\tilde{\varphi}\colon V\times\cdots\times V\to X$ an alternating multilinear map also possessing the universal mapping property for alternating multilinear maps \index{Alternating multilinear maps}

with domain  $ V \times \cdots \times V $, then there is an isomorphism  $ \alpha\colon \Lambda_k(V) \to X $ such that  $ \alpha \circ \varphi = \widetilde{\varphi} $.

In the special case in which  $ W = \mathbb{R} $, the diagram (4) establishes a natural isomorphism

\begin{equation}
\tag{5}\label{eq:2.6:5}
\Lambda_{k}(V)^{*}\cong A_{k}(V)
\end{equation}

of $\Lambda_k(V)^{*}$ with the vector space $A_k(V)$ of all alternating multilinear functions on $V\times\cdots\times V$ (k copies). It follows from property (b) that $A_k(V)=\{0\}$ for $k>\dim V$.

We shall now consider various dualities between the spaces $V_{r,s}$, $\Lambda_{k}(V)$, $\Lambda(V)$ and the corresponding spaces $(V^{*})_{r,s}$, $\Lambda_{k}(V^{*})$, $\Lambda(V^{*})$ built on the dual space $V^{*}$ of $V$.


\warnernumber{7}
\begin{definition}\label{thm:2.7}
Let $V$ and $W$ be real finite dimensional vector spaces. A pairing of
$V$ and $W$ is a bilinear map $(\ ,\ ):V\times W\to\mathbb R$.
A pairing is called non-singular if, for every $w\ne0$ in $W$, there
exists $v\in V$ such that $(v,w)\ne0$, and, for every $v\ne0$ in $V$,
there exists $w\in W$ such that $(v,w)\ne0$. \index{Pairing}

Let $V$ and $W$ be non-singularly paired by $(\ ,\ )$, and define

\begin{equation}
\tag{1}\label{eq:2.7:1}
\varphi\colon V\to W^{*}\quad\text{by}\quad\varphi(v)(w)=(v,w)\qquad(v\in V;~w\in W).
\end{equation}

It follows that $\varphi$ is $1:1$. Similarly, there is a $1:1$ map $W \to V^{*}$. Therefore $V$ and $W$ have the same dimension, and hence $\varphi$ is an isomorphism of $V$ with $W^{*}$. Thus a non-singular pairing of $V$ and $W$ in a canonical way yields an isomorphism $\varphi: V \to W^{*}$ and similarly an isomorphism $W \to V^{*}$.

\end{definition}

\warnernumber{8}
\begin{definition}\label{thm:2.8}
A non-singular pairing of $(V^{*})_{r,s}$ with $V_{r,s}$. This pairing is to be the bilinear map $(V^{*})_{r,s} \times V_{r,s} \to \mathbb{R}$ which on decomposable elements

\[
v^{*}=v_{1}^{*}\otimes\cdots\otimes v_{r}^{*}\otimes u_{r+1}\otimes\cdots\otimes u_{r+s}\in(V^{*})_{r,s}
\]

and

\[
u=u_{1}\otimes\cdots\otimes u_{r}\otimes v_{r+1}^{*}\otimes\cdots\otimes v_{r+s}^{*}\in V_{r,s}
\]

yields the following:

\begin{equation}
\tag{1}\label{eq:2.8:1}
(v^{*},u)=v_{1}^{*}(u_{1})\cdots v_{r+s}^{*}(u_{r+s}).
\end{equation}

It is easily checked that there is a unique such bilinear map and that it is a non-singular pairing. This pairing establishes an isomorphism

\begin{equation}
\tag{2}\label{eq:2.8:2}
(V^{*})_{r,s}\cong(V_{r,s})^{*}.
\end{equation}

On the other hand, the obvious extension of the universal mapping property \hyperref[thm:2.2]{2.2(a)} shows that there is a natural isomorphism

\begin{equation}
\tag{3}\label{eq:2.8:3}
(V_{r,s})^{*}\cong M_{r,s}(V)
\end{equation}

where $M_{r,s}(V)$ is the vector space of all multilinear functions

\[
\underbrace{V\times\cdots\times V}_{r\text{ copies}}\times\underbrace{V^{*}\times\cdots\times V^{*}}_{s\text{ copies}}\to\mathbb{R}.
\]

Under the isomorphism (3), if $\tilde{h} \in (V_{r,s})^{*}$, then the corresponding multi-linear function $h$ in $M_{r,s}(V)$ satisfies

\begin{equation}
\tag{4}\label{eq:2.8:4}
h(v_{1},\ldots,v_{r},v_{1}^{*},\ldots,v_{s}^{*})=\tilde{h}(v_{1}\otimes\cdots\otimes v_{r}\otimes v_{1}^{*}\otimes\cdots\otimes v_{s}^{*}).
\end{equation}

Finally, from (2) and (3) we obtain an isomorphism

\begin{equation}
\tag{5}\label{eq:2.8:5}
(V^{*})_{r,s}\cong M_{r,s}(V).
\end{equation}

\end{definition}

\warnernumber{9}
\begin{definition}\label{thm:2.9}
A non-singular pairing of $\Lambda_k(V^*)$ with $\Lambda_k(V)$. This pairing is to be the bilinear map of $\Lambda_k(V^*) \times \Lambda_k(V) \to \mathbb{R}$ which on decomposable elements $v^* = v_1^* \wedge \cdots \wedge v_k^*$ in $\Lambda_k(V^*)$ and $u = u_1 \wedge \cdots \wedge u_k$ in $\Lambda_k(V)$ yields
\begin{equation}
\tag{1}\label{eq:2.9:1}
(v^{*},u)=\det\big(v_{i}^{*}(u_{j})\big).
\end{equation}


Again it is easily checked that there is a unique such bilinear map and that it is a non-singular pairing. For $k=0$, the pairing is simply multiplication of real numbers. This pairing establishes an isomorphism


\begin{equation}
\tag{2}\label{eq:2.9:2}
\Lambda_{k}(V^{*})\cong\Lambda_{k}(V)^{*}.
\end{equation}

Composing (2) with the natural isomorphism

\begin{equation}
\tag{3}\label{eq:2.9:3}
\Lambda_{k}(V)^{*}\cong A_{k}(V)
\end{equation}

of \eqref{eq:2.6:5} yields an isomorphism

\begin{equation}
\tag{4}\label{eq:2.9:4}
\Lambda_{k}(V^{*})\cong A_{k}(V).
\end{equation}

Using the isomorphism (2), and observing that the dual space of a finite direct sum is canonically isomorphic to the direct sum of the dual spaces, we obtain isomorphism

\begin{equation}
\tag{5}\label{eq:2.9:5}
\Lambda(V^{*})=\sum_{k=0}^{\infty}\Lambda_{k}(V^{*})\cong\sum_{k=0}^{\infty}\Lambda_{k}(V)^{*}\cong\Lambda(V)^{*},
\end{equation}

and from (3) we obtain an isomorphism

\begin{equation}
\tag{6}\label{eq:2.9:6}
\Lambda(V)^{*}\cong A(V)=\sum_{k=0}^{\infty}A_{k}(V).
\end{equation}

Henceforth we shall make use of the identification (5) of $\Lambda(V^{*})$ with $\Lambda(V)^{*}$ via the pairing (1) without further comment.

\end{definition}

\warnernumber{10}
\begin{remarks}{on 2.9}\label{thm:2.10}

(a) If $\{e_1, \ldots, e_d\}$ is a basis of $V$ with dual basis $\{\gamma_1, \ldots, \gamma_d\}$ in $V^*$, then the bases $\{e_\Phi\}$ and $\{\gamma_\Phi\}$ defined in \hyperref[thm:2.6]{2.6(b)} are dual bases of $\Lambda(V)$ and $\Lambda(V^*)$ under the isomorphism \eqref{eq:2.9:5}.

(b) In \eqref{eq:2.9:5} and \eqref{eq:2.9:6}, the following isomorphisms were established:

\[
\Lambda(V^{*})\cong\Lambda(V)^{*}\cong A(V).
\]

The second of these is the natural isomorphism arising from the universal mapping property \hyperref[thm:2.6]{2.6(c)}. The first, call it $\alpha$ for the moment, arose from our choice of the pairing \eqref{eq:2.9:1}. There is another pairing in common use, this one obtained by replacing \eqref{eq:2.9:1} by

\begin{equation}
\tag{1}\label{eq:2.10:1}
(v^{*},u)=\frac{1}{k!}\det\bigl(v_{i}^{*}(u_{j})\bigr).
\end{equation}

This pairing gives a different isomorphism of $\Lambda(V^*)$ with $\Lambda(V)^*,$ call it $\beta$. Now $\Lambda(V^*)$ is an algebra under wedge multiplication. Via the above isomorphism we obtain two algebra structures $\wedge_a$ and $\wedge_\beta$ on $A(V)$. It is not difficult to see that if $f\in A_p(V)$ and $g\in A_q(V)$, then these induced algebra structures on $A(V)$ take the form

\begin{equation}
\tag{2}\label{eq:2.10:2}
\begin{aligned}{f}&{{}\wedge_{\alpha}g(v_{1},\ldots,v_{p+q})}\\ {}&{{}=\sum_{p,q\text{ shuffles}}(\operatorname{sgn}\pi)f(v_{\pi(1)},\ldots,v_{\pi(p)})g(v_{\pi(p+1)},\ldots,v_{\pi(p+q)})}\\ \end{aligned}
\end{equation}

and


\begin{equation}
\tag{3}\label{eq:2.10:3}
\begin{aligned}{f}&{{}\wedge_{\beta}g(v_{1},\ldots,v_{p+q})}\\ {}&{{}=\frac{1}{(p+q)!}\sum_{\pi\in S_{p+q}}(\operatorname{sgn}\pi)f(v_{\pi(1)},\ldots,v_{\pi(p)})g(v_{\pi(p+1)},\ldots,v_{\pi(p+q)}).}\\ \end{aligned}
\end{equation}

Here a permutation  $ \pi \in S_{p+q} $ is called a ``p, q shuffle'' if  $ \pi(1) < \pi(2) < \cdots < \pi(p) $ and  $ \pi(p+1) < \cdots < \pi(p+q) $. It follows from (2) and (3) that

\begin{equation}
\tag{4}\label{eq:2.10:4}
f\wedge_{\alpha}g=\frac{(p+q)!}{p!q!}f\wedge_{\beta}g.
\end{equation}

For example, consider the case in which $p = q = 1$. Let $\gamma$ and $\delta$ belong to $V^* = A_1(V)$, and let $v, w \in V$. Then $\gamma \wedge_\alpha \delta$ and $\gamma \wedge_\beta \delta \in A_2(V)$ and

\begin{equation}
\tag{5}\label{eq:2.10:5}
\gamma\wedge_{\alpha}\delta(v,w)=\gamma(v)\delta(w)-\gamma(w)\delta(v),
\end{equation}

whereas

\begin{equation}
\tag{6}\label{eq:2.10:6}
\gamma\wedge_{\beta}\delta\left(v,w\right)=\tfrac{1}{2}\big(\gamma(v)\,\delta(w)-\gamma(w)\,\delta(v)\big).
\end{equation}

We shall use only the pairing \eqref{eq:2.9:1}. It has the advantage of avoiding factors such as the $\frac{1}{2}$ in (6).

\end{remarks}

\warnernumber{11}
\discuss{2.11}{Linear Transformations of $\Lambda(V)$}\label{thm:2.11}
End$(\Lambda(V))$ will denote the vector space of all endomorphisms of $\Lambda(V)$ (i.e. linear transformations from $\Lambda(V)$ into $\Lambda(V)$). Let $u \in \Lambda(V)$. Left multiplication by $u$ is the endomorphism $e(u)$ of $\Lambda(V)$ defined by \index{Left multiplication}

\begin{equation}
\tag{1}\label{eq:2.11:1}
\varepsilon(u)v=u\land v\qquad(v\in\Lambda(V)).
\end{equation}

The transpose $i(u)$ of $\varepsilon(u)$ is an endomorphism of $\Lambda(V)^{*}$. Under our identification of $\Lambda(V)^{*}$ with $\Lambda(V^{*})$, the transpose $i(u)$ can be considered as an endomorphism of $\Lambda(V^{*})$, and this endomorphism is called interior multiplication by $u$. In terms of our pairing of $\Lambda(V)$ with $\Lambda(V^{*})$, $i(u)$ is defined by \index{Interior multiplication}

\begin{equation}
\tag{2}\label{eq:2.11:2}
\big(i(u)v^{*},w\big)=\big(v^{*},\varepsilon(u)w\big)=(v^{*},u\land w)\qquad\big(v^{*}\in\Lambda(V^{*});w\in\Lambda(V)\big).
\end{equation}

If $u \in V$, then $i(u)$ maps $\Lambda_k(V^*)$ into $\Lambda_{k-1}(V^*)$ for each $k$. In particular, if $v^* \in V^*$, then $i(u)v^* \in \mathbb{R}$ and

\[
i(u)v^{*}=\left(i(u)v^{*},1\right)=\left(v^{*},\varepsilon(u)\cdot1\right)=(v^{*},u)=v^{*}(u).
\]

An endomorphism $l$ of $\Lambda(V)$ (or in general of any graded algebra) is

(a) a derivation if $l(u \land v) = l(u) \land v + u \land l(v)$ $\quad (u, v \in \Lambda(V))$,

(b) an anti-derivation if $l(u \land v) = l(u) \land v + (-1)^p u \land l(v)$

$(u \in \Lambda_p(V); v \in \Lambda(V))$,

(c) of degree k if  $ l\colon \Lambda_{j}(V) \to \Lambda_{j+k}(V) $ for all j.

(We assume that  $ \Lambda_{i}(V) = \{0\} $ if  $ i < 0 $.)

It is easy to see that  $ l \in \text{End} \Lambda(V) $ is an anti-derivation if and only if on decomposable elements we have

\begin{equation}
\tag{3}\label{eq:2.11:3}
l(v_{1}\wedge\cdots\wedge v_{j})=\sum_{i=1}^{j}(-1)^{i+1}v_{1}\wedge\cdots\wedge l(v_{i})\wedge\cdots\wedge v_{j}.
\end{equation}


\warnernumber{12}
\begin{proposition}\label{thm:2.12}
If $u \in V$, then $i(u)$ is an anti-derivation of degree -1. \index{Anti-derivation}

\end{proposition}

\begin{proof}
That the degree of $i(u)$ is $-1$ if $u \in V$ is clear from the \hyperref[thm:2.11]{definition 2.11}(2). To prove that $i(u)$ is an anti-derivation, we check that \eqref{eq:2.11:3} holds. For this, it is sufficient to see that both sides of \eqref{eq:2.11:3} give the same result when paired with an element of $\Lambda(V)$ of the form $w_{2} \wedge \cdots \wedge w_{j}$. That is, we must show that

\begin{equation}
\tag{1}\label{eq:2.12:1}
\begin{aligned}{}&{{}\left(i(u)(v_{1}^{*}\wedge\cdots\wedge v_{j}^{*}),w_{2}\wedge\cdots\wedge w_{j}\right)}\\ {}&{{}\qquad=\left(\sum_{i=1}^{j}(-1)^{i+1}v_{1}^{*}\wedge\cdots\wedge i(u)v_{i}^{*}\wedge\cdots\wedge v_{j}^{*},w_{2}\wedge\cdots\wedge w_{j}\right).}\\ \end{aligned}
\end{equation}



Now, the left-hand side of (1) equals

\[
(v_{1}^{*}\wedge\cdots\wedge v_{j}^{*},u\wedge w_{2}\wedge\cdots\wedge w_{j})=\det(v_{i}^{*}(w_{l})),
\]

where we have put  $ w_{1}=u $. The right-hand side of (1) is


\begin{align*}\sum_{i=1}^{j}(-1)^{i+1}v_{i}^{*}(u)&(v_{1}^{*}\wedge\cdots\wedge\widehat{v_{i}^{*}}\wedge\cdots\wedge v_{j}^{*},w_{2}\wedge\cdots\wedge w_{j})\\&=\sum_{i=1}^{j}(-1)^{i+1}v_{i}^{*}(u)\det\big(v_{k}^{*}(w_{l})\big)\quad\begin{pmatrix}k=1,\ldots,\widehat{i},\ldots,j\\ l=2,\ldots,j\end{pmatrix}\\&=\det\big(v_{i}^{*}(w_{l})\big)\quad\begin{pmatrix}i=1,\ldots,j\\ l=1,\ldots,j\end{pmatrix}.\end{align*}


(Here the circumflex over a term means that it is to be omitted.)

\end{proof}

\warnernumber{13}
\discuss{2.13}{Effect of Linear Transformations}\label{thm:2.13}
Let $l: V \to W$ be a linear transformation. Then $l$ extends to an algebra homomorphism


\begin{equation}
\tag{1}\label{eq:2.13:1}
l\colon\Lambda(V)\to\Lambda(W),
\end{equation}

where

\[
l(v_{1}\wedge\cdots\wedge v_{k})=l(v_{1})\wedge\cdots\wedge l(v_{k})\quad\mathrm{and}\quad l(1)=1.
\]

The transpose  $ \delta l $:  $ W^{*}\rightarrow V^{*} $ also extends to an algebra homomorphism

\begin{equation}
\tag{2}\label{eq:2.13:2}
\delta l\colon\Lambda(W^{*})\to\Lambda(V^{*}).
\end{equation}

Our isomorphisms \eqref{eq:2.9:5} of  $ \Lambda(V^{*}) $ with  $ \Lambda(V)^{*} $ and of  $ \Lambda(W^{*}) $ with  $ \Lambda(W)^{*} $ are natural in the sense that (2) is the transpose of (1); that is,

\begin{equation}
\tag{3}\label{eq:2.13:3}
\big(\delta l(w^{*}),v\big)=\big(w^{*},l(v)\big)\qquad\big(w^{*}\in\Lambda(W^{*});v\in\Lambda(V)\big).
\end{equation}

\section{Tensor Fields and Differential Forms}


\warnernumber{14}
\begin{definition}\label{thm:2.14}
Let M be a differentiable manifold. We define \index{Differential forms}

\[
T_{r,s}(M)=\bigcup_{m\in M}(M_{m})_{r,s}
\]

tensor bundle of type (r,s) over M

\[
\Lambda_{k}^{*}(M)=\bigcup_{m\in M}\Lambda_{k}(M_{m}^{*})
\]

exterior k bundle over M


\[
\Lambda^{*}(M)=\bigcup_{m\in M}\Lambda(M_{m}^{*})
\]

exterior algebra bundle over M. \index{Exterior k-bundle}\index{Exterior algebra bundle}

In the cases in which $k = 0$ and $(r,s) = (0,0)$, the unions in (1) and (2) are disjoint unions of copies of the real line—one copy for each point in $M$. $T_{r,s}(M)$, $\Lambda_{k}^{*}(M)$, and $\Lambda^{*}(M)$ have natural manifold structures such that the canonical projection maps to $M$ are $C^{\infty}$. If $(U,\varphi)$ is a coordinate system on $M$ with coordinate functions $y_{1}, \ldots, y_{d}$, then the bases $\{\partial/\partial y_{i}\}$ of $M_{m}$ and $\{dy_{i}\}$ of $M_{m}^{*}$, for $m \in U$, yield bases of $(M_{m})_{r,s}$, $\Lambda_{k}(M_{m}^{*})$, and $\Lambda(M_{m}^{*})$. For example, the basis of $\Lambda_{k}(M_{m}^{*})$ is $\{dy_{i_{1}} \wedge \cdots \wedge dy_{i_{k}}: i_{1} < \cdots < i_{k}\}$.

Using these bases, one can define maps of the inverse images of $U$ in $T_{r,s}(M)$, $\Lambda_{k}^{*}(M)$, and $\Lambda^{*}(M)$ under the respective projection maps to $\varphi(U) \times$ Euclidean spaces of the proper dimensions. By requiring these maps to be coordinate systems, one obtains the natural manifold structures on $T_{r,s}(M)$, $\Lambda_{k}^{*}(M)$, and $\Lambda^{*}(M)$ just as we did previously in 1.25 for $T(M)$ and $T^{*}(M)$ (which incidentally is simply $\Lambda_{1}^{*}(M)$).

\end{definition}

\warnernumber{15}
\begin{definition}\label{thm:2.15}
A $C^{\infty}$ mapping of $M$ into $T_{r,s}(M)$, $\Lambda_{k}^{*}(M)$, or $\Lambda^{*}(M)$ whose composition with the canonical projection is the identity map is called a (smooth) tensor field of type $(r,s)$ on $M$, a (differential) $k$-form on $M$, or a (differential) form on $M$ respectively. Since our tensor fields and differential forms will always be smooth in this sense, we shall drop the adjectives “smooth” and “differential” unless needed for emphasis. \index{Tensor fields}

\end{definition}

\warnernumber{16}
\begin{remarks}\label{thm:2.16}
A lifting  $ \alpha\colon M \to T_{r,s}(M) $ is a smooth tensor field of type  $ (r,s) $ if and only if for each coordinate system  $ (U, y_1, \ldots, y_d) $ on  $ M $,

\begin{equation}
\tag{1}\label{eq:2.16:1}
\alpha\mid U=\sum a_{i_{1},\ldots,i_{r};j_{1},\ldots,j_{s}}\frac{\partial}{\partial y_{i_{1}}}\otimes\cdots\otimes\frac{\partial}{\partial y_{i_{r}}}\otimes d y_{j_{1}}\otimes\cdots\otimes d y_{j_{s}},
\end{equation}

where the  $ a_{i_1,\ldots,i_r;j_1,\ldots,j_s} \in C^\infty(U) $. (In the classical tensor notation one uses lower indices on tangent vectors, upper indices on functions and differentials, and the reverse on coefficients; thus the individual terms of (1) would appear as

\[
a_{j_{1},\ldots,j_{s}}^{i_{1},\ldots,i_{r}}\frac{\partial}{\partial y^{i_{1}}}\otimes\cdots\otimes\frac{\partial}{\partial y^{i_{r}}}\otimes d y^{j_{1}}\otimes\cdots\otimes d y^{j_{s}}.)\nonumber
\]

A lifting $\beta\colon M\to\Lambda_{k}^{*}(M)$ is a differential $k$-form if and only if for each coordinate system $(U,y_{1},\ldots,y_{d})$ on $M$,

\begin{equation}
\tag{2}\label{eq:2.16:2}
\beta\mid U=\sum_{i_{1}<\dots<i_{k}}b_{i_{1},\dots,i_{k}}d y_{i_{1}}\wedge\dots\wedge d y_{i_{k}},
\end{equation}

where the  $ b_{i_{1},...,i_{k}} $ are  $ C^{\infty} $ functions on U.

\end{remarks}

\warnernumber{17}
\begin{definition}\label{thm:2.17}
$ E^k(M) $ shall denote the set of all smooth k-forms on M, and  $ E^*(M) $ the set of all differential forms.  $ E^0(M) $ can be identified with  $ C^\infty(M) $; indeed, the differentiable manifold  $ \Lambda_0^*(M) $ is simply  $ M \times \mathbb{R} $, and smooth liftings of M into  $ M \times \mathbb{R} $ are simply graphs of  $ C^\infty $ functions on M. Forms can be added, multiplied by scalars, and given a product  $ (\wedge) $. If  $ \omega, \varphi \in E^*(M) $ and  $ c \in \mathbb{R} $, then  $ \omega + \varphi $,  $ c\omega $, and  $ \omega \wedge \varphi $ are the forms which at m have the values  $ \omega_m + \varphi_m $,  $ c\omega_m $, and  $ \omega_m \wedge \varphi_m $ respectively. In the case in which f is a 0-form and  $ \omega \in E^*(M) $, we write  $ f \wedge \omega $ simply as  $ f\omega $.  $ E^*(M) $ thus has the structure both of a module over the ring  $ C^\infty(M) $ and of a graded algebra over  $ \mathbb{R} $ with wedge multiplication.

\end{definition}

\warnernumber{18}
\discuss{2.18}{}\label{thm:2.18}
Let  $ \omega \in E^k(M) $. Then  $ \omega_m \in \Lambda_k(M_m^*) $, and can be considered (via the duality \eqref{eq:2.9:4}) as an alternating multilinear function on  $ M_m $. So if  $ X_1, \ldots, X_k $ are vector fields on  $ M $,  $ \omega(X_1, \ldots, X_k) $ makes sense—it is the function whose value at  $ m $ is

\begin{equation}
\tag{1}\label{eq:2.18:1}
\omega(X_{1},\ldots,X_{k})(m)=\omega_{m}\bigl(X_{1}(m),\ldots,X_{k}(m)\bigr).
\end{equation}

Thus, if we let  $ \mathfrak{X}(M) $ denote the  $ C^\infty(M) $ module of smooth vector fields on M, then

\begin{equation}
\tag{2}\label{eq:2.18:2}
\omega\colon\underbrace{\mathfrak{X}(M)\times\cdots\times\mathfrak{X}(M)}_{k\text{ copies}}\to C^{\infty}(M)
\end{equation}

and is an alternating multilinear map of the module  $ \mathfrak{X}(M) $ into  $ C^\infty(M) $. We stress that  $ \omega $ is multilinear over the  $ C^\infty(M) $ module  $ \mathfrak{X}(M) $; that is

\begin{equation}
\tag{3}\label{eq:2.18:3}
\begin{aligned}{\omega(X_{1},\ldots,X_{i-1},}&{{}f X+g Y,X_{i+1},\ldots,X_{k})}\\ {\quad=}&{{}f\omega(X_{1},\ldots,X_{i-1},X,X_{i+1},\ldots,X_{k})}\\ {}&{{}+g\omega(X_{1},\ldots,X_{i-1},Y,X_{i+1},\ldots,X_{k})}\\ \end{aligned}
\end{equation}

whenever $f, g \in C^\infty(M)$ and $X_1, \ldots, X_{i-1}, X, Y, X_{i+1}, \ldots, X_k \in \mathfrak{X}(M)$.

Conversely, it is useful to observe that any alternating $C^\infty(M)$ multilinear map (2) of the module $\mathfrak{X}(M)$ into $C^\infty(M)$ defines a form; for we claim that if $\omega$ is such a map, then $\omega(X_1, \ldots, X_k)(m)$ depends only on the values of the vector fields $X_i$ at $m$. Assuming this for the moment, it follows that $\omega$ defines an alternating multilinear function $\omega_m$ on $M_m$, and hence defines an element of $\Lambda_k(M_m^*)$; namely, given $(v_1, \ldots, v_k) \in M_m \times \cdots \times M_m$, choose $V_1, \ldots, V_k \in \mathfrak{X}(M)$ such that $V_i(m) = v_i$ ($i = 1, \ldots, k$), and define

\begin{equation}
\tag{4}\label{eq:2.18:4}
\omega_{m}(v_{1},\ldots,v_{k})=\omega(V_{1},\ldots,V_{k})(m).
\end{equation}

By our claim,  $ \omega_m(v_1, \ldots, v_k) $ is well-defined, independent of the choice of the extensions  $ V_i $. Thus  $ \omega $ gives a lifting  $ m \mapsto \omega_m $ of  $ M $ into  $ \Lambda^*(M) $, and this is easily seen to be smooth; hence  $ \omega $ is a form.

For simplicity of notation, we illustrate the claim with the case in which $\omega$ is a linear map of the module $\mathfrak{X}(M)$ into $C^\infty(M)$. Let $X\in\mathfrak{X}(M)$. We wish to show that $\omega(X)(m)$ depends only on $X(m)$. It suffices to show that $\omega(X)(m)=0$ if $X(m)=0$. Let $(U,x_1,\ldots,x_d)$ be a coordinate system about $m$. Then on $U$, we have $X=\sum a_i(\partial/\partial x_i)$ where $a_i(m)=0$. Now, let $\varphi$ be a $C^\infty$ function which takes the value 1 on a neighborhood $V\subset U$ of $m$ and is zero on a neighborhood of $M-U$ (see 1.10). Then the vector field $X_i$ which is $\varphi(\partial/\partial x_i)$ on $U$ and $0$ elsewhere is a $C^\infty$ vector field on $M$; the function $\bar{a}_t$ which is $\varphi a_t$ on $U$ and $0$ elsewhere belongs to $C^\infty(M)$; and

\begin{equation}
\tag{5}\label{eq:2.18:5}
X=\sum\tilde{a}_{i}X_{i}+(1-\varphi^{2})X.
\end{equation}

Thus

\begin{equation}
\tag{6}\label{eq:2.18:6}
\begin{array}{r}{\omega(X)(m)=\sum\tilde{a}_{i}(m)\omega(X_{i})(m)+\big((1-\varphi^{2})(m)\big)\big(\omega(X)(m)\big)=0.}\end{array}
\end{equation}

So  $ \omega(X)(m)=0 $ if  $ X(m)=0 $, as was to be shown.

Finally, we remark that via the duality \eqref{eq:2.8:5}, tensor fields can be given a similar interpretation. If $T$ is a tensor field of type $(r,s)$, then we can consider $T$ as a map

\begin{equation}
\tag{7}\label{eq:2.18:7}
T\colon\underbrace{E^{1}(M)\times\cdots\times E^{1}(M)}_{r\text{ copies}}\times\underbrace{\mathfrak{X}(M)\times\cdots\times\mathfrak{X}(M)}_{s\text{ copies}}\to C^{\infty}(M),
\end{equation}

which is $C^{\infty}(M)$ multilinear with respect to the $C^{\infty}(M)$ modules $E^{1}(M)$ and $\mathfrak{X}(M)$.

Observe the simple form which formula (2) of \hyperref[thm:2.10]{2.10(b)} takes when $\omega, \varphi \in E^1(M)$ and $X, Y \in \mathfrak{X}(M)$. In this case, we have \index{Exterior derivative}

\begin{equation}
\tag{8}\label{eq:2.18:8}
\omega\wedge\varphi(X,Y)=\omega(X)\varphi(Y)-\omega(Y)\varphi(X).
\end{equation}


\warnernumber{19}
\begin{definition}\label{thm:2.19}
If $f \in C^\infty(M)$, the differential $df$ is a smooth mapping of $T(M)$ into $\mathbb{R}$ which is linear on each tangent space. Thus $df$ can be considered as a 1-form, $df: M \to \Lambda_1^*(M)$. The 1-form $df$ is called the exterior derivative of the 0-form $f$, and this exterior differentiation operator $d$ has an important extension to $E^*(M)$ given by the following.

\end{definition}

\warnernumber{20}
\begin{theorem}{Exterior Differentiation}\label{thm:2.20}
There exists a unique anti-derivation $d\colon E^{*}(M)\to E^{*}(M)$ of degree +1 such that

(1)  $ d^{2}=0 $

(2) Whenever  $ f \in C^\infty(M) = E^0(M) $, df is the differential of f.

\end{theorem}

\begin{proof}
Existence. Let $p \in M$. Let $E^*(p)$ be the set of all smooth forms defined on open subsets of $M$ containing $p$, with $E^k(p)$ the corresponding set of $k$-forms. We fix a coordinate system $(U, x_1, \ldots, x_d)$ about $p$. If $\omega \in E^*(p)$, then


\[
\omega\big|_{(\text{domain }\omega)\cap U}=\sum a_{\Phi}d x_{\Phi}
\]

where the $a_\Phi \in C^\infty\big((\text{domain} \omega) \cap U\big)$, where $\Phi$ runs over all subsets of $\{1, \ldots, d\}$, and where the $dx_\Phi$ are either $dx_{\epsilon_1} \wedge \cdots \wedge dx_{\epsilon_r}$ when $\Phi = \{i_1 < \cdots < i_r\}$ or the constant function $1$ when $\Phi = \varnothing$. We define $d\omega$ at $p$ by setting

\[
\begin{array}{r}{d\omega_{p}=\sum d a_{\Phi}|_{p}\wedge d x_{\Phi}|_{p}\in\Lambda(M_{p}^{*}).}\end{array}
\]


We will have to show that the definition of  $ d\omega $, is independent of the choice of coordinates. But first, we give the following properties:

(a)  $ \omega \in E^{r}(p) \Rightarrow d\omega_{p} \in \Lambda_{r+1}(M_{p}^{*}) $.

(b)  $ d\omega_{p} $ depends only on the germ of  $ \omega $ at p.

(c)  $ d(a_1\omega_1 + a_2\omega_2)|_{p} = a_1(d\omega_1)|_{p} + a_2(d\omega_2)|_{p} $,  $ (a_i \in \mathbb{R}; \omega_i \in E^*(p)) $, where the domain of  $ a_1\omega_1 + a_2\omega_2 $ is (domain $ \omega_1\ \cap $ domain $ \omega_2 $).

\[
\begin{array}{r}{d(\omega_{1}\wedge\omega_{2})\big|_{p}=d\omega_{1}\big|_{p}\wedge\omega_{2}\big|_{p}+(-1)^{r}\omega_{1}\big|_{p}\wedge d\omega_{2}\big|_{p}}\\ {\big(\omega_{1}\in E^{r}(p);\omega_{2}\in E^{*}(p)\big).}\end{array}
\]

In view of properties (b) and (c), it suffices to check (d) for $\omega_1 = f dx_{i_1} \wedge \cdots \wedge dx_{i_r}$ and $\omega_2 = g\, dx_{j_1} \wedge \cdots \wedge dx_{j_s}$ on some neighborhood of $p$. For the case $r = s = 0$, property (d) is simply $d(f \cdot g)|_{p} = df_p \cdot g(p) + f(p) \cdot dg_p$; and the case in which only one of $r$ or $s$ is $0$ is similar. Now suppose that $r > 0$ and $s > 0$. If $\{i_1, \ldots, i_r\} \cap \{j_1, \ldots, j_s\} \neq \varnothing$, both sides are 0. So assume that this intersection is empty. Then

\[
\begin{split}(f\, dx_{i_{1}}\wedge\cdots\wedge dx_{i_{r}})\wedge(g\, dx_{j_{1}}\wedge\cdots\wedge dx_{j_{s}})&\\&=\varepsilon f\cdot g\, dx_{l_{1}}\wedge\cdots\wedge dx_{l_{r+s}}\end{split}
\]

where  $ l_{1} < \cdots < l_{r+s} $ and  $ \varepsilon $ is the sign of the permutation that has been carried out. So we have that

\begin{align*}d(\omega_{1}\wedge\omega_{2})\big|_{p}&=d\big(\varepsilon f\cdot g\, dx_{l_{1}}\wedge\cdots\wedge dx_{l_{r+s}}\big)\big|_{p}\\&=\varepsilon\big(df_{p}\cdot g(p)+f(p)\, dg_{p}\big)\wedge dx_{l_{1}}\big|_{p}\wedge\cdots\wedge dx_{l_{r+s}}\big|_{p}\\&=\big(df_{p}\wedge dx_{i_{1}}\big|_{p}\wedge\cdots\wedge dx_{i_{r}}\big|_{p}\big)\wedge\big(g(p)\, dx_{j_{1}}\big|_{p}\wedge\cdots\wedge dx_{j_{s}}\big|_{p}\big)\\&\qquad+(-1)^{r}\big(f(p)\, dx_{i_{1}}\big|_{p}\wedge\cdots\wedge dx_{i_{r}}\big|_{p}\big)\\&\qquad\qquad\wedge\big(dg_{p}\wedge dx_{j_{1}}\big|_{p}\wedge\cdots\wedge dx_{j_{s}}\big|_{p}\big)\\&=d\omega_{1}\big|_{p}\wedge\omega_{2}\big|_{p}+(-1)^{r}\omega_{1}\big|_{p}\wedge d\omega_{2}\big|_{p}.\end{align*}

(e) If $f$ is a $C^\infty$ function on a neighborhood of $p$, then $d(df)|_{p}=0$. For on (domain $f)\cap U$, $df=\sum(\partial f/\partial x_i)dx_i$, so that

\[
d(d f)|_{p}=\sum d\left(\frac{\partial f}{\partial x_{i}}\right)\Bigg|_{p}\wedge d x_{i}\big|_{p}=\sum_{i,j}\frac{\partial^{2}f}{\partial x_{j}\partial x_{i}}\Bigg|_{p}d x_{j}\big|_{p}\wedge d x_{i}\big|_{p}.
\]

But  $ (\partial^2f/\partial x_j\,\partial x_i)(p)=(\partial^2f/\partial x_i\,\partial x_j)(p) $, whereas  $ dx_j\big|_{p}\wedge dx_i\big|_{p}=-dx_i\big|_{p}\wedge dx_j\big|_{p} $. So  $ d(df)\big|_{p}=0 $.

Now we claim that the definition of $d$ at $p$ is independent of the coordinates chosen. For let $d'$ be defined on $E^*(p)$ relative to another coordinate system, and let $\omega \in E^*(p)$. Then $\omega$ on (domain $\omega$) $\cap U$ is given by (3), and since $d'$ must also satisfy properties (a)–(e), it follows that

\[
\begin{array}{r l r}{d^{\prime}(\omega)\big\vert_{p}}&{=d^{\prime}(\sum a_{\Phi}d x_{i_{1}}\wedge\dots\wedge d x_{i_{r}})\big\vert_{p}}&{\quad\mathrm{(by~(b))}}\\ &{=\sum d^{\prime}(a_{\Phi}d x_{i_{1}}\wedge\dots\wedge d x_{i_{r}})\big\vert_{p}}&{\quad\mathrm{(by~(c))}}\\ &{=\sum d^{\prime}(a_{\Phi})\big\vert_{p}\wedge d x_{i_{1}}\big\vert_{p}\wedge\dots\wedge d x_{i_{r}}\big\vert_{p}}\\ &{\quad+\sum(-1)^{k-1}a_{\Phi}\big\vert_{p}d x_{i_{1}}\big\vert_{p}\wedge\dots\wedge d^{\prime}(d x_{i_{k}})\big\vert_{p}\wedge\dots\wedge d x_{i_{r}}\big\vert_{p}}\\ &{}&{\quad\mathrm{(by~(d))}}\\ &{=\sum d(a_{\Phi})\big\vert_{p}\wedge d x_{i_{1}}\big\vert_{p}\wedge\dots\wedge d x_{i_{r}}\big\vert_{p}}&{\quad\mathrm{(by~(e))}}\\ &{=d\omega_{p}.}\end{array}
\]

Now, if $\omega \in E^*(M)$, we define $d\omega$ to be the form which as a lifting of $M$ into $\Lambda^*(M)$ sends $p$ to $d\omega_p$. It follows that $d^2 = 0$, since if $\omega \in E^*(M)$ and $p \in M$, $d\omega$ has the form $\sum da_\Phi \wedge dx_\Phi$ on a coordinate neighborhood of $p$, and thus

\[
d(d\omega)|_{p}=\sum d(d a_{\Phi}\wedge d x_{\Phi})|_{p}=0
\]

by properties (e) and (d). It follows that $d$ is an anti-derivation of $E^{*}(M)$ of degree +1 satisfying (1) and (2).

Uniqueness. Let $d'$ also be an anti-derivation of $E^{*}(M)$ of degree +1 satisfying (1) and (2). We first show that if $\omega \in E^{*}(M)$ and $\omega$ vanishes on a neighborhood $W$ of $p$, then $d'\omega|_{p} = 0$. Choose a $C^{\infty}$ function $\varphi$ which is 0 on a neighborhood of $p$ and 1 on a neighborhood of $M - W$. Then $\varphi\omega = \omega$ and

\[
d^{\prime}(\omega)\big|_{p}=d^{\prime}(\varphi\omega)\big|_{p}=d^{\prime}(\varphi)\big|_{p}\wedge\omega_{p}+\varphi(p)\,d^{\prime}\omega\big|_{p}=0.
\]

Now $d'$ is defined only on elements of $E^{*}(M)$, that is, on globally defined forms on $M$. We wish to define $d'$ on $E^{*}(p)$ for each $p \in M$. If $\omega \in E^{*}(p)$, we can extend $\omega$ to a form on $M$ having the same germ at $p$ as does $\omega$. Simply let $\varphi$ be a $C^{\infty}$ function which is 1 on a neighborhood of $p$ and has support in the domain of $\omega$. Then $\varphi\omega \in E^{*}(M)$ ($\varphi\omega$ is defined to be 0 outside of the domain of $\omega$), and $\varphi\omega$ agrees with $\omega$ on a neighborhood of $p$. Thus we can define

\[
d^{\prime}(\omega)|_{p}=d^{\prime}(\varphi\omega)|_{p},
\]

and by the above remarks this definition is independent of the extension chosen. Thus  $ d'(\omega)|_{p} $ is defined for all  $ \omega $ in  $ E^{*}(p) $ and clearly satisfies properties (a)–(e). (In (d), extend  $ \omega_{1} \wedge \omega_{2} $ to a form on M by using  $ \varphi\omega_{1} \wedge \varphi\omega_{2} $ for a suitable  $ \varphi $; and in (e), observe that  $ d'(df)|_{p}=d'(\varphi df)|_{p}=d'(d(\varphi f))|_{p}=0 $, since  $ d\varphi(p)=0 $ and since  $ d' $ satisfies (1) and (2).) The equations (5) now imply that whenever  $ \omega \in E^{*}(p) $ in particular whenever  $ \omega \in E^{*}(M) $,  $ d'(\omega)|_{p}=d(\omega)|_{p} $. This proves uniqueness.

Observe that it is clear from the above proof that  $ d\omega \mid U = d(\omega \mid U) $ whenever U is an open set in M.

\end{proof}

\warnernumber{21}
\discuss{2.21}{Interior Multiplication by Vector Fields}\label{thm:2.21}
Let $X$ be a smooth vector field on $M$, and let $\omega \in E^*(M)$. Interior multiplication of $\omega$ by $X$ is the form $i(X)\omega$ whose value at $m$ is the interior multiple of $\omega_m$ by $X_m$ (see \eqref{eq:2.11:2}):

\begin{equation}
\tag{1}\label{eq:2.21:1}
\left(i(X)\omega\right)\big|_{m}=i(X_{m})(\omega_{m}).
\end{equation}

That $i(X)\omega$ is smooth follows easily from \eqref{eq:2.16:2}; and from 2.12 it follows that $i(X)\colon E^{*}(M)\to E^{*}(M)$ is an anti-derivation of degree -1.

\warnernumber{22}
\discuss{2.22}{Effect of Mappings}\label{thm:2.22}
Let $\psi: M \to N$ be a smooth map, and let $m \in M$. Then we have the differential $d\psi: M_m \to N_{\psi(m)}$, its transpose $\delta\psi: N_{\psi(m)}^* \to M_m^*$, and the induced algebra homomorphism $\delta\psi\colon \Lambda\big(N_{\psi(m)}^{*}\big) \to \Lambda(M_{m}^{*})$. If $\omega$ is a form on $N$, then we can pull $\omega$ back to a form on $M$ by setting

\begin{equation}
\tag{1}\label{eq:2.22:1}
\delta\psi(\omega)|_{m}=\delta\psi(\omega|_{\psi(m)}).
\end{equation}

This is one of the particularly nice features of differential forms. Under a smooth mapping, they can be pulled back from the range to the domain of the map. Vector fields, on the other hand, do not display such pleasing behavior under mappings.


\warnernumber{23}
\begin{proposition}\label{thm:2.23}
Let  $ \psi\colon M \to N $ be a smooth map. Then

(a)  $ \delta\psi\colon E^{*}(N)\to E^{*}(M) $ and is an algebra homomorphism.

(b)  $ \delta\psi $ commutes with  $ d $; that is,  $ d(\delta\psi(\omega)) = \delta\psi(d\omega) $ ( $ \omega \in E^*(N) $).

(c)  $ \delta\psi(\omega)(X_1,\ldots,X_k)(m)=\omega_{\psi(m)}(d\psi(X_1(m)),\ldots,d\psi(X_k(m))) $ for  $ \omega\in E^k(N) $ and for vector fields  $ X_1,\ldots,X_k $ on  $ M $.

\end{proposition}

\begin{proof}
Result (c) is clear from \eqref{eq:2.13:3}; and \eqref{eq:2.13:2} and the \hyperref[thm:2.22]{definition 2.22}(1) imply that $\delta\psi$ is an algebra homomorphism. Before we check that $\delta\psi(\omega)$ is actually a smooth form for $\omega\in E^{*}(N)$, observe the following special case of (b): If $f\in C^\infty(N)$, then $\delta\psi(f)=f\circ\psi\in C^\infty(M)$, and \hyperref[thm:1.23]{1.23(d)} implies that

\begin{equation}
\tag{1}\label{eq:2.23:1}
\delta\psi(d f)=d(f\circ\psi)=d(\delta\psi(f)).
\end{equation}

Now let $\omega \in E^{*}(N)$, and let $m \in M$. Choose a coordinate system $(U, x_1, \ldots, x_d)$ about $\psi(m)$ and a neighborhood $V$ of $m$ such that $\psi(V) \subset U$. Then there are $C^\infty$ functions $a_\Phi$ on $U$ such that

\begin{equation}
\tag{2}\label{eq:2.23:2}
\omega\mid U=\sum a_{\Phi}d x_{i_{1}}\wedge\cdots\wedge d x_{i_{r}}.
\end{equation}

It follows that

\begin{equation}
\tag{3}\label{eq:2.23:3}
\delta\psi(\omega)\bigm|V=\sum a_{\Phi}\circ\psi\;d(x_{i_{1}}\circ\psi)\wedge\dots\wedge d(x_{i_{r}}\circ\psi),
\end{equation}

which is a smooth form on $V$. Hence $\delta\psi(\omega)\in E^{*}(M)$, and thus (a) is proved. To complete (b), we use (3):

\begin{equation}
\tag{4}\label{eq:2.23:4}
\begin{aligned}{d\big(\delta\psi(\omega)\big)\big|_{m}}&{{}=d\big(\textstyle\sum a_{\Phi}\circ\psi d(x_{i_{1}}\circ\psi)\wedge\cdots\wedge d(x_{i_{r}}\circ\psi)\big)\big|_{m}}\\ {}&{{}=\textstyle\sum\big(d(a_{\Phi}\circ\psi)\wedge d(x_{i_{1}}\circ\psi)\wedge\cdots\wedge d(x_{i_{r}}\circ\psi)\big)\big|_{m}}\\ {}&{{}=\delta\psi(\textstyle\sum d a_{\Phi}\wedge d x_{i_{1}}\wedge\cdots\wedge d x_{i_{r}}\big)\big|_{m}}\\ {}&{{}=\delta\psi(d\omega)\big|_{m}.}\\ \end{aligned}
\end{equation}

Thus (b) is proved, and the proof is complete.

\end{proof}

\section{The Lie Derivative}


\warnernumber{24}
\begin{definition}\label{thm:2.24}
Tensor fields and differential forms can be differentiated with respect to a vector field. The resulting derivative is known as the Lie derivative and is defined as follows. Fix a smooth vector field $X$ on a manifold $M$. Recall (see 1.48) that we use $X_t$ to denote the local 1-parameter group of transformations associated with $X$. Let $Y$ be another smooth vector field on $M$. We shall define the derivative of $Y$ with respect to $X$ at the point $m \in M$. First we follow the integral curve of $X$ through $m$ out to the point $X_t(m)$ and evaluate $Y$ there. Then we transfer $Y_{X_{t}(m)}$ back to $M_m$ via the differential $dX_{-t}$ of the diffeomorphism $X_{-t}$. In $M_m$ we take the difference of the vectors $dX_{-t}(Y_{X_{t}(m)})$ and $Y_m$, divide the difference by $t$, and then take the limit as $t \to 0$. In other words, we consider the smooth $M_m$-valued function $t \mapsto dX_{-t}(Y_{X_{t}(m)})$, and we take its derivative at $t = 0$. The result is a vector in $M_m$ which is called the Lie derivative of $Y$ with respect to $X$ at $m$ and which is denoted by $(L_X Y)_m$. Thus we define

\begin{equation}
\tag{1}\label{eq:2.24:1}
(L_{X}Y)_{m}=\lim_{t\to0}\frac{d X_{-t}(Y_{X_{t}(m)})-Y_{m}}{t}=\frac{d}{d t}\bigg|_{t=0}\big(d X_{-t}(Y_{X_{t}(m)})\big).
\end{equation}

\begin{center}
\redrawn{figure-24}
\end{center}

In a similar way we define the Lie derivative of a differential form $\omega$ with respect to the vector field $X$, except that in this case we evaluate $\omega$ at $X_t(m)$ and then pull back via $\delta X_t$ to $\Lambda(M_m^*)$ where we take the difference with $\omega(m)$, divide by $t$, and take the limit as $t \to 0$. Thus we define \index{Lie derivative}

\begin{equation}
\tag{2}\label{eq:2.24:2}
(L_{X}\omega)_{m}=\lim_{t\to0}\frac{\delta X_{t}(\omega_{X_{t}(m)})-\omega_{m}}{t}=\frac{d}{d t}\bigg|_{t=0}\big(\delta X_{t}(\omega_{X_{t}(m)})\big).
\end{equation}

Smoothness of $L_{X}\omega$ and of $L_{X}Y$ is asserted in \hyperref[thm:2.25]{Proposition 2.25}. The Lie derivative $L_{X}$ can be extended to arbitrary tensor fields in the obvious way. If $T$ is a tensor field of type $(r,s)$, then $(L_{X}T)_{m}$ is the derivative at $t=0$ of the $(M_{m})_{r,s}$-valued function whose value at $t$ is

\[
d X_{-t}(v_{1}\otimes\cdots\otimes v_{r})\otimes\delta X_{t}(v_{1}^{*}\otimes\cdots\otimes v_{s}^{*})
\]

if

\[
T\big|_{X_{t}(m)}=v_{1}\otimes\cdots\otimes v_{r}\otimes v_{1}^{*}\otimes\cdots\otimes v_{s}^{*}.
\]

\end{definition}

\warnernumber{25}
\begin{proposition}\label{thm:2.25}
Let $X$ be a $C^{\infty}$ vector field on $M$. Then

(a)  $ L_{X}f = X(f) $ whenever  $ f \in C^{\infty}(M) $.

(b)  $ L_{X}Y = [X,Y] $ for each  $ C^{\infty} $ vector field Y on M.

(c)  $ L_{X}\colon E^{*}(M)\to E^{*}(M) $, and is a derivation which commutes with  $ d $.

(d) On  $ E^{*}(M) $,  $ L_{X} = i(X) \circ d + d \circ i(X) $.

(e) Let $\omega \in E^{p}(M)$, and let $Y_{0}, \ldots, Y_{p}$ be $C^{\infty}$ vector fields on $M$. Then

\[
\begin{aligned}{L_{Y_{0}}\big(\omega(Y_{1},\ldots,Y_{p})\big)}&{{}=(L_{Y_{0}}\omega)(Y_{1},\ldots,Y_{p})}\\ {}&{{}+\sum_{i=1}^{p}\omega(Y_{1},\ldots,Y_{i-1},L_{Y_{0}}Y_{i},Y_{i+1},\ldots,Y_{p}).}\\ \end{aligned}
\]

(f) Assumption as in (e). Then

\[
\begin{aligned}{d\omega(Y_{0},\ldots,Y_{p})}&{{}=\sum_{i=0}^{p}(-1)^{i}Y_{i}\omega(Y_{0},\ldots,\widehat{Y_{i}},\ldots,Y_{p})}\\ {+}&{{}\sum_{i<j}(-1)^{i+j}\omega([Y_{i},Y_{j}],\;Y_{0},\ldots,\widehat{Y_{i}},\ldots,\widehat{Y_{j}},\ldots,Y_{p}).}\\ \end{aligned}
\]

\end{proposition}

\begin{proof}
We leave result (a) as an exercise. For (b), we need only show that  $ L_{X}Y(f)=[X,Y] $ for each  $ f\in C^{\infty}(M) $. Let  $ m\in M $. Then

\begin{equation}
\tag{1}\label{eq:2.25:1}
\begin{aligned}{(L_{X}Y)_{m}(f)}&{{}=\bigg(\lim_{t\to0}\frac{d X_{-t}Y_{X_{t}(m)}-Y_{m}}{t}\bigg)(f)}\\ {}&{{}=\frac{d}{d t}\bigg|_{t=0}\big[\big(d X_{-t}(Y_{X_{t}(m)})\big)(f)\big]}\\ {}&{{}=\frac{d}{d t}\bigg|_{t=0}[Y_{X_{t}(m)}(f\circ X_{-t})].}\\ \end{aligned}
\end{equation}

Define a real-valued function $H$ on a neighborhood of $(0,0)$ in $\mathbb{R}^{2}$ by setting

\begin{equation}
\tag{2}\label{eq:2.25:2}
H(t,u)=f\bigl(X_{-t}(Y_{u}(X_{t}(m)))\bigr).
\end{equation}

Then by (a),

\begin{equation}
\tag{3}\label{eq:2.25:3}
Y_{X_{t}(m)}(f\circ X_{-t})=\frac{\partial}{\partial r_{2}}\bigg\vert_{(t,0)}H(t,u),
\end{equation}

and (1) becomes

\begin{equation}
\tag{4}\label{eq:2.25:4}
(L_{X}Y)_{m}(f)=\frac{\partial^{2}H}{\partial r_{1}\partial r_{2}}\bigg|_{(0,0)}.
\end{equation}

To evaluate this derivative, we set

\begin{equation}
\tag{5}\label{eq:2.25:5}
K(t,u,s)=f\bigl(X_{s}(Y_{u}(X_{t}(m)))\bigr).
\end{equation}

Then  $ H(t,u) = K(t,u,-t) $. It follows from the chain rule that

\begin{equation}
\tag{6}\label{eq:2.25:6}
\frac{\partial^{2}H}{\partial r_{1}\partial r_{2}}\bigg|_{(0,0)}=\frac{\partial^{2}K}{\partial r_{1}\partial r_{2}}\bigg|_{(0,0,0)}-\frac{\partial^{2}K}{\partial r_{3}\partial r_{2}}\bigg|_{(0,0,0)}.
\end{equation}

Now,  $ K(t,u,0) = f(Y_u(X_t(m))) $. Hence

\begin{equation}
\tag{7}\label{eq:2.25:7}\begin{gathered}
\left.\frac{\partial K}{\partial r_{2}}\right|_{(t,0,0)}=Y_{X_{t}(m)}(f),\\[4pt]
\frac{\partial^{2}K}{\partial r_{1}\partial r_{2}}\bigg|_{(0,0,0)}=X_{m}(Y f).
\end{gathered}
\end{equation}

Also  $ K(0,u,s) = f(X_s(Y_u(m))) $. Hence

\begin{equation}
\tag{8}\label{eq:2.25:8}\begin{gathered}
\frac{\partial K}{\partial r_{3}}\bigg|_{(0,u,0)}=X f\big(Y_{u}(m)\big),\\[4pt]
\frac{\partial^{2}K}{\partial r_{2}\partial r_{3}}\bigg|_{(0,0,0)}=Y_{m}(X f).
\end{gathered}
\end{equation}

Part (b) now follows from (4), (6), (7), and (8). One immediate consequence of (b) is that $L_{X}Y$ is a smooth vector field.

For result (c), first consider $L_{X}$ as a mapping of $E^{*}(M)$ into forms which are not a priori smooth, and observe that the derivation property follows immediately from adding and subtracting suitable terms before taking the limit in the definition \eqref{eq:2.24:1}. Next we check that $L_{X}$ commutes with $d$ when applied to functions; that is,

\begin{equation}
\tag{9}\label{eq:2.25:9}
\big(L_{X}(d f)\big)_{m}=d(L_{X}f)_{m}\qquad\big(f\in C^{\infty}(M);m\in M\big).
\end{equation}

Since both sides of (9) are elements of  $ M_{m}^{*} $, we need only prove that they have the same effect when applied to an arbitrary vector  $ Y_{m} $ in  $ M_{m} $. The right-hand side gives

\begin{equation}
\tag{10}\label{eq:2.25:10}
d(L_{X}f)_{m}(Y_{m})=Y_{m}(L_{X}f)=Y_{m}\biggl(\frac{d}{d t}\bigg\vert_{t=0}(f\circ X_{t})\biggr),
\end{equation}

where $f\circ X_{t}$ can be considered as a $C^{\infty}$ function on $(-\varepsilon,\varepsilon)\times W$ for some $\varepsilon>0$ and some neighborhood $W$ of $m$ in $M$ (1.48(d)). The left-hand side of (9) gives

\begin{equation}
\tag{11}\label{eq:2.25:11}
\begin{aligned}{\big(L_{X}(d f)\big)_{m}(Y_{m})}&{{}=\bigg\{\frac{d}{d t}\bigg|_{t=0}\big(\delta X_{t}(d f_{X_{t}(m)})\big)\bigg\}(Y_{m})}\\ {}&{{}=\frac{d}{d t}\bigg|_{t=0}\big(\delta X_{t}(d f_{X_{t}(m)})(Y_{m})\big)}\\ {}&{{}=\frac{d}{d t}\bigg|_{t=0}\big(d f(d X_{t}(Y_{m}))\big)=\frac{d}{d t}\bigg|_{t=0}\big(Y_{m}(f\circ X_{t})\big).}\end{aligned}.
\end{equation}

Now let $Y$ be an extension of $Y_{m}$ to a vector field on $W$. Then, according to Exercise 24(c) of Chapter 1, $d/dt$ and $Y$ have canonical extensions to vector fields on $(-\varepsilon,\varepsilon)\times W$ where they satisfy $[d/dt,Y] \equiv 0$. This fact together with (10) and (11) implies (9). Finally, to see that the form $L_{X}\omega$ is actually smooth and to check that $L_{X}$ commutes with $d$ on all of $E^{*}(M)$, simply express an arbitrary form $\omega$ in local coordinates as in \eqref{eq:2.16:2}, and compute using equation (9), result (a), and the fact that $L_{X}$ is a derivation.

For result (d), observe that both $L_{X}$ and $i(X)\circ d+d\circ i(X)$ are derivations on $E^{*}(M)$ which commute with $d$, and both have the same effect on functions. Then (d) follows from a simple computation in local coordinates.

We shall leave (e) as an exercise. The proof is not difficult—one simply has to be persistent enough in unwinding the definitions. Again, one checks the identity by restricting to a local coordinate system. Try it first for the case  $ \omega = fdx_1 \wedge dx_2 $ in a coordinate neighborhood.

Finally, result (f) follows from (e) by using (d) and induction on p. For the case p = 1, we have from (e) that

\begin{equation}
\tag{12}\label{eq:2.25:12}
L_{Y_{0}}\big(\omega(Y_{1})\big)=(L_{Y_{0}}\omega)(Y_{1})+\omega([Y_{0},Y_{1}]).
\end{equation}

Applying (a) and (d) to (12), one obtains

\[
\begin{aligned}{Y_{0}\omega(Y_{1})}&{{}=\big((i(Y_{0})\circ d+d\circ i(Y_{0}))\omega\big)(Y_{1})+\omega([Y_{0},Y_{1}])}\\ {}&{{}=d\omega(Y_{0},Y_{1})+Y_{1}\big(\omega(Y_{0})\big)+\omega([Y_{0},Y_{1}]),}\\ \end{aligned}
\]

which is result (f) in the case p = 1. Now assume that (f) holds for p - 1. Then (f) is obtained for p by again starting with (e) and applying (d) and the induction hypothesis.

\end{proof}

\section{Differential Ideals}


Our objective here is to give a version of the Frobenius theorem in terms of differential forms, and then to describe E. Cartan's useful method of obtaining maps by looking for their graphs.

\warnernumber{26}
\begin{definition}\label{thm:2.26}
Let $\mathcal{D}$ be a $p$-dimensional $C^{\infty}$ distribution on $M$. A $q$-form $\omega$ is said to annihilate $\mathcal{D}$ if for each $m\in M$

$ \omega_m(v_1, \ldots, v_q) = 0 $ whenever  $ v_1, \ldots, v_q \in \mathcal{D}(m) $.

A form  $ \omega \in E^*(M) $ is said to annihilate  $ \mathcal{D} $ if each of the homogeneous parts of  $ \omega $ annihilates  $ \mathcal{D} $. We let

$ \mathcal{I}(\mathcal{D}) = \{\omega \in E^*(M)\colon \omega \text{ annihilates } \mathcal{D}\}. $

\end{definition}

\warnernumber{27}
\begin{definition}\label{thm:2.27}
A collection  $ \omega_1, \ldots, \omega_n $ of 1-forms on  $ M $ is called independent if they form an independent set in  $ M_m^* $ for each  $ m \in M $.

\end{definition}

\warnernumber{28}
\begin{proposition}\label{thm:2.28}
Let $\mathcal{D}$ be a smooth $p$-dimensional distribution on M. Then

(a)  $ \mathcal{I}(\mathcal{D}) $ is an ideal in  $ E^{*}(M) $.

(b) $\mathcal{I}(\mathcal{D})$ is locally generated by $d-p$ independent 1-forms. (That is, to each $m\in M$ there corresponds a neighborhood $U$ of $m$ and a set of independent 1-forms $\omega_{1},\ldots,\omega_{d-p}$ on $U$ such that:

(i) If $\omega \in \mathcal{I}(\mathcal{D})$, then $\omega \mid U$ belongs to the ideal in $E^{*}(U)$ generated by $\omega_1, \ldots, \omega_{d-p}$.

(ii) If $\omega \in E^*(M)$, and if there is a cover of $M$ by sets $U$ (as above) such that for each $U$ in the cover, $\omega \mid U$ belongs to the ideal generated by $\omega_1, \ldots, \omega_{d-p}$, then $\omega \in \mathcal{I}(\mathcal{D})$.)

(c) If $\mathcal{I} \subset E^{*}(M)$ is an ideal locally generated by $d - p$ independent 1-forms, then there exists a unique $C^{\infty}$ distribution $\mathcal{D}$ of dimension $p$ on $M$ for which $\mathcal{I} = \mathcal{I}(\mathcal{D})$.

\end{proposition}

\begin{proof}
Part (a) follows from the definition of  $ \mathcal{I}(\mathcal{D}) $ and the definition of multiplication in  $ E^{*}(M) $.

Let $m \in M$. Since $\mathcal{D}$ is smooth and $p$-dimensional, there exist $C^\infty$ vector fields $X_{d-p+1},\ldots,X_d$ defined and spanning $\mathcal{D}$ at each point of a neighborhood of $m$. This collection can be completed to a collection $X_1,\ldots,X_d$ of smooth vector fields forming a basis of $M_n$ for each $n$ in a neighborhood $U$ of $m$. Let $\omega_1,\ldots,\omega_d$ be the dual 1-forms; that is,

\[
\omega_{i}(X_{j})(n)=\delta_{i j}\qquad\text{(Kronecker index)}
\]

for each $n$ in $U$. Then $\omega_{1},\ldots,\omega_{d-p}$ are the desired 1-forms on $U$.



They are independent smooth 1-forms on $U$. If $\omega \in \mathcal{I}(\mathcal{D})$, $\omega \mid U = \sum a_{\Phi}\, \omega_{i_{1}} \wedge \cdots \wedge \omega_{i_{r}}$, where $\Phi$ runs over nonempty subsets $\{i_1, \ldots, i_r\} \subset \{1, \ldots, d\}$, and where the $a_{\Phi}$ must be identically zero unless $\{i_1, \ldots, i_r\} \cap \{1, \ldots, d-p\} \neq \emptyset$. Thus $\omega \mid U$ belongs to the ideal in $E^*(U)$ generated by $\omega_{1}, \ldots, \omega_{d-p}$. Conversely, if $\omega$ is a form such that for each such $U$ in some covering of $M$, $\omega \mid U$ belongs to the ideal in $E^*(U)$ generated by $\omega_1, \ldots, \omega_{d-p}$, then clearly $\omega \in \mathcal{I}(\mathcal{D})$. This proves (b).

For part (c), let $m \in M$, and let the independent 1-forms $\omega_1, \ldots, \omega_{d-p}$ generate $\mathcal{I}$ on a neighborhood $U$ of $m$. Define $\mathcal{D}(m)$ to be the subspace of $M_m$ whose annihilator is the subspace of $M_m^*$ spanned by the collection $\{\omega_{i}(m): i=1, \ldots, d-p\}$. It follows that $\mathcal{D}$ is a smooth $p$-dimensional distribution on $M$ and that $\mathcal{I} = \mathcal{I}(\mathcal{D})$. Uniqueness of $\mathcal{D}$ follows from the fact that $\mathcal{D} \neq \mathcal{D}_1$ implies $\mathcal{I}(\mathcal{D}) \neq \mathcal{I}(\mathcal{D}_1)$.
\end{proof}

\warnernumber{29}
\begin{definition}\label{thm:2.29}
An ideal $\mathcal{I} \subset E^{*}(M)$ is called a differential ideal if it is closed under exterior differentiation $d$; that is,

\begin{equation}
\tag{1}\label{eq:2.29:1}
d(\mathcal{I})\subset\mathcal{I}.
\end{equation}

\end{definition}

\warnernumber{30}
\begin{proposition}\label{thm:2.30}
A $C^{\infty}$ distribution $\mathcal{D}$ on $M$ is involutive if and only if the ideal $\mathcal{I}(\mathcal{D})$ is a differential ideal.

\end{proposition}

\begin{proof}
Let $\omega$ be a $q$-form in $\mathcal{I}(\mathcal{D})$, and let $X_0, \ldots, X_q$ be smooth vector fields lying in $\mathcal{D}$. Then $2.25(\mathrm{f})$ together with the involutiveness of $\mathcal{D}$ implies that $d\omega(X_0, \ldots, X_q) \equiv 0$. Hence $d\omega \in \mathcal{I}(\mathcal{D})$, and $\mathcal{I}(\mathcal{D})$ is a differential ideal. Conversely, suppose that $\mathcal{I}(\mathcal{D})$ is a differential ideal. Let $Y_0$ and $Y_1$ be vector fields lying in $\mathcal{D}$, and let $m \in M$. By $2.28(\mathrm{b})$, there are independent 1-forms $\omega_1, \ldots, \omega_{d-p}$ generating $\mathcal{I}(\mathcal{D})$ on a neighborhood $U$ of $m$. Extend these forms to $M$ by multiplying by a $C^\infty$ function which is 1 on a neighborhood of $m$ and has support in $U$. We shall denote the extended forms similarly by $\omega_1, \ldots, \omega_{d-p}$. By $2.25(\mathrm{f})$,

\begin{equation}
\tag{1}\label{eq:2.30:1}
\omega_{i}([Y_{0},Y_{1}])=-d\omega_{i}(Y_{0},Y_{1})+Y_{0}\omega_{i}(Y_{1})-Y_{1}\omega_{i}(Y_{0})
\end{equation}

for $i=1,\ldots,d-p$. The right-hand side of (1) is identically zero on $M$ since $\mathcal{I}(\mathcal{D})$ is a differential ideal and since $\omega_i\in\mathcal{I}(\mathcal{D})$. Thus $\omega_i([Y_0,Y_1])(m)=0$ for $i=1,\ldots,d-p$. Now $\mathcal{D}(m)$ is the subspace of $M_m$ whose annihilator is the subspace of $M_m^*$ spanned by the collection $\{\omega_i(m):i=1,\ldots,d-p\}$. Thus $[Y_0,Y_1](m)\in\mathcal{D}(m)$, and $\mathcal{D}$ is involutive.

\end{proof}

\warnernumber{31}
\begin{definition}\label{thm:2.31}
A submanifold $(N,\psi)$ of $M$ is an integral manifold of an ideal $\mathcal{I} \subset E^{*}(M)$ if for every $\omega \in \mathcal{I}$, $\delta\psi(\omega) \equiv 0$. A connected integral manifold of an ideal $\mathcal{I}$ is maximal if its image is not a proper subset of the image of any other connected integral manifold of the ideal.

It now follows immediately that we have the following version of the Frobenius \hyperref[thm:1.64]{theorem 1.64} in terms of differential ideals.

\end{definition}

\warnernumber{32}
\begin{theorem}\label{thm:2.32}
Let $\mathcal{I} \subset E^{*}(M)$ be a differential ideal locally generated by $d-p$ independent 1-forms. Let $m \in M$. Then there exists a unique maximal, connected, integral manifold of $\mathcal{I}$ through $m$, and this integral manifold has dimension $p$.

\end{theorem}

\warnernumber{33}
\discuss{2.33}{}\label{thm:2.33} We shall now consider a technique which in certain situations enables one to find a map by looking for its graph as an integral manifold of a differential ideal. This will have important applications in the theory of Lie groups in Chapter 3, and also has important applications in such areas as isometric imbeddings in Riemannian geometry.

Suppose that $f\colon N^{c}\to M^{d}$ is $C^{\infty}$ and that $\{\omega_{i}\}$ is some collection of forms on $M$. Let $\pi_{1}$ and $\pi_{2}$ denote the natural projections of $N\times M$ onto $N$ and $M$ respectively. For each $i$, we define a form $\mu_{i}$ on $N\times M$ by setting

\begin{equation}
\tag{1}\label{eq:2.33:1}
\mu_{i}=\delta\pi_{1}\delta f(\omega_{i})-\delta\pi_{2}(\omega_{i}).
\end{equation}

Let  $ \mathcal{I} $ be the ideal in  $ E^{*}(N \times M) $ generated by the  $ \mu_i $.

Now the graph of f is the submanifold $(N,g)$ of $N\times M$ where \index{Graph of a map}

\begin{equation}
\tag{2}\label{eq:2.33:2}
g(n)=(n,f(n)).
\end{equation}

We claim that the graph is an integral manifold of the ideal $\mathcal{I}$. For this, it suffices to show that $\delta g(\mu_{i})=0$ for each $i$. Now, $\pi_{1}\circ g=\mathrm{id}$, and $\pi_{2}\circ g=f$, and thus it follows that

\[
\delta g(\mu_{i})=\delta(\pi_{1}\circ g)\delta f(\omega_{i})-\delta(\pi_{2}\circ g)(\omega_{i})=\delta f(\omega_{i})-\delta f(\omega_{i})=0.
\]

Thus starting with a map $f\colon N\to M$ and a collection of forms on $M$, we have observed that the graph of $f$ is an integral manifold of a certain ideal of forms on $N\times M$. Now, suppose that we start with the manifold $M^d$, and suppose that there exists a basis $\omega_1,\ldots,\omega_d$ of 1-forms on $M$, that is, $\{\omega_i(m)\}$ is a basis of $M_m^*$ for each $m\in M$. (Such a basis does not generally exist, but it does exist in a number of interesting situations for which the following technique is quite useful.) Suppose also that we have a manifold $N^c$ and a collection $\alpha_1,\ldots,\alpha_d$ of 1-forms on $N$ and that we wish to find a map $f\colon N\to M$ such that

\begin{equation}
\tag{3}\label{eq:2.33:3}
\delta f(\omega_{i})=\alpha_{i}
\end{equation}

for $i=1,\ldots,d$. If the map $f$ exists, then, as we have observed, its graph will be an integral manifold of a certain ideal of forms. So we attempt to find the graph from the ideal. We define forms $\mu_{i}$ on $N\times M$ by setting

\begin{equation}
\tag{4}\label{eq:2.33:4}
\mu_{i}=\delta\pi_{1}(\alpha_{i})-\delta\pi_{2}(\omega_{i}),
\end{equation}

and we let $\mathcal{I}$ be the ideal in $E^{*}(N \times M)$ which they generate. If $\mathcal{I}$ happens to be a differential ideal, then we can obtain the desired map $f$ (at least locally) from an integral manifold of $\mathcal{I}$. That is, one gets at $f$ by looking for

its graph as an integral manifold of a suitable differential ideal. So suppose that $\mathcal{I}$ is a differential ideal, and let $(n_{0},m_{0})\in N\times M$. Then since $\mathcal{I}$ is locally (in fact globally) generated by $d$ independent 1-forms, the Frobenius theorem guarantees that there is a maximal, connected, integral manifold $I$ of $\mathcal{I}$ of dimension $c$ through $(n_{0},m_{0})$. Let $q\in I$. We claim that $d\pi_{1}\mid I_{q}$ is 1:1. For suppose that $v\in I_{q}$ and that $d\pi_{1}(v)=0$. Then since $\mu_{i}(v)=0$, it follows from (4) that $\omega_{i}(d\pi_{2}(v))=0$ for $i=1,\ldots,d$; and this implies that $d\pi_{2}(v)=0$ since the $\omega_{i}$ form a basis of 1-forms on $M$. But now if both $d\pi_{1}(v)=0$ and $d\pi_{2}(v)=0$, then $v=0$. Hence $d\pi_{1}\mid I_{q}$ is 1:1. Thus $\pi_{1}\mid I:I\to N$ is locally a diffeomorphism. So there exist neighborhoods $V$ of $(n_{0},m_{0})$ in $I$ and $U$ of $n_{0}$ such that $\pi_{1}\mid V:V\to U$ is a diffeomorphism. We define $f:U\to M$ by setting

\begin{equation}
\tag{5}\label{eq:2.33:5}
f=\pi_{2}\circ(\pi_{1}\mid V)^{-1}.
\end{equation}

Then  $ f(n_{0}) = m_{0} $, the graph of f is an open submanifold of $I$, and moreover,

\begin{equation}
\tag{6}\label{eq:2.33:6}
\delta f(\omega_{i})=\alpha_{i}\mid U.
\end{equation}

For let  $ v \in U_{n} $ for some  $ n \in U $. Then by (4),

\[
\begin{aligned}{0}&{{}=\mu_{i}\big(d(\pi_{1}\mid V)^{-1}(v)\big)}\\ {}&{{}=\alpha_{i}(v)-\omega_{i}\big(d\pi_{2}\circ d(\pi_{1}\mid V)^{-1}(v)\big)=\alpha_{i}(v)-\delta f(\omega_{i})(v),}\\ \end{aligned}
\]

which proves (6). Thus we have obtained the desired map locally. That is, given $n_{0} \in N$ and an arbitrary choice of $m_{0} \in M$, then under the assumption that the ideal $\mathcal{I}$ generated by (4) is a differential ideal, we have found an open neighborhood $U$ of $n_{0}$ and a $C^{\infty}$ map $f: U \to M$ such that $f(n_{0}) = m_{0}$ and such that $\delta f(\omega_{i}) = \alpha_{i} \mid U$. Moreover, there is a unique such map. More precisely, if $m_{0} \in M$, and if $U$ is any connected open neighborhood of $n_{0}$ in $N$ for which there exists a $C^{\infty}$ map $f: U \to M$ such that $f(n_{0}) = m_{0}$ and such that $\delta f(\omega_{i}) = \alpha_{i} \mid U$ for $t = 1, \ldots, d$, then there is a unique such map on $U$. For let $\tilde{f}$ be any other such map. Let $(U, \tilde{g})$ and $(U, g)$ be the graphs of $\tilde{f}$ and $f$ over $U$ respectively. Thus

\begin{equation}
\tag{7}\label{eq:2.33:7}
\tilde{g}(n)=\left(n,\tilde{f}(n)\right)\quad\text{and}\quad g(n)=\left(n,f(n)\right)
\end{equation}

for $n \in U$. Then, according to our remarks near the beginning of this section, not only is $(U,g)$ an integral manifold of $\mathcal{I}$ through $(n_{0},m_{0})$, but so is $(U,\tilde{g})$. Now, the subset of $U$ on which $g$ and $\tilde{g}$ agree is non-empty since it contains $n_{0}$, and is closed by continuity, and is moreover open. For let $\tilde{g}(n) = g(n)$. Then it follows from the uniqueness of integral manifolds that there exist sufficiently small neighborhoods $W$ and $\tilde{W}$ of $n$ so that

\begin{equation}
\tag{8}\label{eq:2.33:8}
g(W)=\tilde{g}(\tilde{W}).
\end{equation}

It follows then from (7) that  $ W = \tilde{W} $ and that

\begin{equation}
\tag{9}\label{eq:2.33:9}
g\mid W=\tilde{g}\mid W.
\end{equation}

Hence, since $U$ is connected, $g = \tilde{g}$ on $U$, which implies that $f = \tilde{f}$ on $U$. This proves uniqueness.

In suitable situations one can conclude that  $ \pi_{1}|I $ is a covering of  $ N $; and then if  $ N $ is simply connected, one can conclude that there exists a unique  $ C^{\infty} $ map  $ f: N \to M $ such that  $ f(n_{0}) = m_{0} $ and such that (3) holds. In Chapter 3 we shall make several concrete applications of this method for proving existence and uniqueness of certain maps. This technique was originally used by E. Cartan for the problem of finding local isometric imbeddings of Riemannian manifolds in Euclidean space.

We summarize the results of this section in the following.


\warnernumber{34}
\begin{theorem}\label{thm:2.34}
Let $N^{c}$ and $M^{d}$ be differentiable manifolds, and let $\pi_{1}$ and $\pi_{2}$ be the canonical projections of $N\times M$ onto $N$ and $M$ respectively. Suppose that there exists a basis $\{\omega_{i}: i=1,\ldots,d\}$ for the 1-forms on $M$.

(a) If $f: N \to M$ is $C^{\infty}$, then the graph of $f$ is an integral manifold of the ideal of forms on $N \times M$ generated by

\begin{equation}
\tag{1}\label{eq:2.34:1}
\{\delta\pi_{1}\delta f(\omega_{i})-\delta\pi_{2}(\omega_{i})\colon i=1,\ldots,d\}.
\end{equation}

(b) If $\{\alpha_{i}\colon i=1,\ldots,d\}$ are $1$-forms on $N$, and if the ideal of forms on  $ N \times M $ generated by

\begin{equation}
\tag{2}\label{eq:2.34:2}
\{\delta\pi_{1}(\alpha_{i})-\delta\pi_{2}(\omega_{i})\colon i=1,\ldots,d\}
\end{equation}

is a differential ideal, then given $n_{0}\in N$ and $m_{0}\in M$ there exists a neighborhood $U$ of $n_{0}$ and a $C^{\infty}$ map $f\colon U\to M$ such that $f(n_{0})=m_{0}$ and such that


\[
\delta f(\omega_{i})=\alpha_{i}\big|_{U}\quad(i=1,\ldots,d).
\]

Moreover, if $U$ is any connected open set containing $n_{0}$ for which there exists a $C^{\infty}$ map $f: U \to M$ satisfying both $f(n_{0}) = m_{0}$ and equation (3), then there exists a unique such map on $U$.

\end{theorem}

\section{Exercises}

\begin{problemset}

\item Supply the proofs of \hyperref[thm:2.2]{2.2(a)}–(e).

\item (a) Show that homogeneous tensors are generally not decomposable.

(b) Show that if $\dim V \leq 3$, then every homogeneous element in $\Lambda(V)$ is decomposable.

(c) Let $\dim V > 3$. Give an example of an indecomposable homogeneous element of $\Lambda(V)$.

(d) Let  $ \alpha $ be a differential form. Is  $ \alpha \wedge \alpha \equiv 0 $?

\item Supply the proofs of \hyperref[thm:2.6]{2.6(a)}–(c).

\item Derive the formulas (2), (3), and (4) of 2.10.

\item Prove that  $ L_{X}f = Xf $ whenever  $ f \in C^{\infty}(M) $ and X is a  $ C^{\infty} $ vector field on M.

\item Let $X$ and $Y$ be $C^\infty$ vector fields on $M$ with corresponding local 1-parameter groups $X_t$ and $Y_t$. Let $m \in M$, and let

\[
\beta(t)=Y_{-\sqrt{t}}\big(X_{-\sqrt{t}}\big(Y_{\sqrt{t}}\big(X_{\sqrt{t}}(m)\big)\big)\big)
\]

for $t$ in $(-\varepsilon,\varepsilon)$, for a sufficiently small $\varepsilon$.

Prove that

\begin{center}
\redrawn{figure-25}
\end{center}

\begin{equation}
\left[X,Y\right]\big|_{m}f=\lim_{t\rightarrow0}\frac{f\big(\beta(t)\big)-f\big(\beta(0)\big)}{t}.
\end{equation}

If $\beta(t)$ were a smooth curve at $t=0$, then the right-hand side of (1) would simply be the effect of the tangent vector to $\beta$ at $t=0$ applied to the function $f$. However, because of the $\sqrt{t}$, $\beta$ is not generally smooth at $t=0$. Thus the existence of the limit in (1) is part of the problem, and the problem asserts that this curve $\beta$, even though it is not smooth at $t=0$, defines a tangent vector in $M_m$ in the usual way, and this vector is precisely $[X,Y]|_{m}$.

\item Prove \hyperref[thm:2.25]{2.25(e)}.

\item Let $M$ be connected, and let $\pi: M \times N \to N$ be the natural projection. Prove that a $p$-form $\omega$ on $M \times N$ is $\delta\pi(\alpha)$ for some $p$-form $\alpha$ on $N$ if and only if $i(X)\omega = 0$ and $L_{X}\omega = 0$ for every vector field $X$ on $M \times N$ for which $d\pi(X(m,n)) = 0$ at each point $(m,n) \in M \times N$.

\item Prove that the elements  $ v_{1}, \ldots, v_{r} $ of the vector space V are linearly independent if and only if  $ v_{1} \wedge \cdots \wedge v_{r} \neq 0 $.

\item Prove that linearly independent sets $\{v_1, \ldots, v_r\}$ and $\{w_1, \ldots, w_r\}$ are bases of the same $r$-dimensional subspace of a vector space $V$ if and only if $v_1 \wedge \cdots \wedge v_r = cw_1 \wedge \cdots \wedge w_r$ where necessarily $c \neq 0$; and in this case, $c = \det A$ where $A = (a_{ij})$ and $v_i = \sum a_{ij}w_j$.

\item Let $\mathcal{I}$ be an ideal of forms on $M$ locally generated by $r$ independent 1-forms. Say $\mathcal{I}$ is generated by $\omega_{1},\ldots,\omega_{r}$ on $U$. Then the condition that $\mathcal{I}$ be a differential ideal is equivalent to each of:

(a)  $ d\omega_i = \sum_j \omega_{ij} \wedge \omega_j $ for some 1-forms  $ \omega_{ij} $ (for each such  $ (U, \omega_1, \ldots, \omega_r) $).

(b) If  $ \omega = \omega_1 \wedge \cdots \wedge \omega_r $, then  $ d\omega = \alpha \wedge \omega $ for some 1-form  $ \alpha $ (for each such  $ (U, \omega_1, \ldots, \omega_r) $).

\item Let $V$ be an $n$-dimensional vector space, and let $l$ be a linear transformation on $V$. Since $\Lambda_n(V)$ is one-dimensional, the linear transformation which $l$ induces on $\Lambda_n(V)$ is simply multiplication by a constant. Define the determinant of $l$ to be this constant. If $A$ is an $n \times n$ matrix, let $v_1, \ldots, v_n$ be a basis of $V$, and let $l$ be the linear transformation on $V$ whose matrix with respect to this basis is $A$. Then define the determinant of $A$ to be the determinant of $l$. Prove that the determinant of $A$ does not depend on the basis $v_1, \ldots, v_n$ chosen. Using this definition, derive the standard properties of determinants. For example, derive the expansion

\[
\det A=\sum_{\pi}(\operatorname{sgn}\pi)a_{1\pi(1)}\cdots a_{n\pi(n)}
\]

where  $ A=(a_{ij}) $,  $ \mathrm{sgn}\pi $ is the sign of the permutation  $ \pi $, and  $ \pi $ runs over all permutations on n letters. Prove also that the determinant of the product of two matrices is the product of their determinants. \index{Determinants}

\item Let $V$ be an $n$-dimensional real inner product space. We extend the inner product from $V$ to all of $\Lambda(V)$ by setting the inner product of elements which are homogeneous of different degrees equal to zero, and by setting

\begin{equation}
\langle w_{1}\wedge\cdots\wedge w_{p},v_{1}\wedge\cdots\wedge v_{p}\rangle=\det\langle w_{i},v_{j}\rangle
\end{equation}

and then extending bilinearly to all of $\Lambda_{p}(V)$. Prove that if $e_{1},\ldots,e_{n}$ is an orthonormal basis of $V$, then the corresponding basis \eqref{eq:2.6:1} of $\Lambda(V)$ is an orthonormal basis for $\Lambda(V)$.

Since $\Lambda_{n}(V)$ is one-dimensional, $\Lambda_{n}(V) - \{0\}$ has two components. An orientation on $V$ is a choice of a component of $\Lambda_{n}(V) - \{0\}$. If $V$ is an oriented inner product space, then there is a linear transformation

\begin{equation}
*:\Lambda(V)\to\Lambda(V),
\end{equation}

called \emph{star}, which is well-defined by the requirement that for any orthonormal basis $e_{1},\ldots,e_{n}$ of $V$ (in particular, for any re-ordering of a given basis), \index{Star operator}


\begin{equation}\begin{aligned}*(1)&=\pm e_{1}\wedge\cdots\wedge e_{n},\qquad*(e_{1}\wedge\cdots\wedge e_{n})=\pm1,\\&\quad*(e_{1}\wedge\cdots\wedge e_{p})=\pm e_{p+1}\wedge\cdots\wedge e_{n},\end{aligned}\end{equation}


where one takes “+” if  $ e_1 \wedge \cdots \wedge e_n $ lies in the component of  $ \Lambda_n(V) - \{0\} $ determined by the orientation and “−” otherwise.

Observe that

\begin{equation}
\ast\colon\Lambda_{p}(V)\to\Lambda_{n-p}(V).
\end{equation}

Prove that on  $ \Lambda_{p}(V) $,

\begin{equation}
**=(-1)^{p(n-p)}.
\end{equation}

Also prove that for arbitrary $v, w \in \Lambda_{p}(V)$, their inner product is given by


\[
\langle v,w\rangle=\ast(w\wedge\ast v)=\ast(v\wedge\ast w).
\]

\item Let $V$ be a real inner product space, as in Exercise 13. Let $\gamma: \Lambda_{p+1}(V) \to \Lambda_p(V)$ be the adjoint of left exterior multiplication by $\xi \in V$. That is,

\[
\langle\gamma(v),w\rangle=\langle v,\xi\wedge w\rangle
\]

for  $ v \in \Lambda_{p+1}(V) $ and  $ w \in \Lambda_p(V) $. Prove that

\[
\gamma(v)=(-1)^{np}*\big(\xi\wedge(*v)\big).
\]

\item Let  $ \xi \in V $. Prove that the composition

\[
\Lambda_{p}(V)\xrightarrow{\xi\land}\Lambda_{p+1}(V)\xrightarrow{\xi\land}\Lambda_{p+2}(V)
\]

of left exterior multiplication by $\xi$ with itself is an exact sequence; that is, the image of the first map is the kernel of the second.

\item Cartan Lemma Let $p \leq d$, and let $\omega_{1}, \ldots, \omega_{p}$ be 1-forms on $M^{d}$ which are linearly independent pointwise. Let $\theta_{1}, \ldots, \theta_{p}$ be 1-forms on $M$ such that

\[
\sum_{i=1}^{p}\theta_{i}\wedge\omega_{i}=0.
\]

Prove that there exist $C^{\infty}$ functions $A_{ij}$ on $M$ with $A_{ij}=A_{ji}$ such that

\[
\theta_{i}=\sum_{j=1}^{p}A_{i j}\omega_{j}\qquad(i=1,\ldots,p).
\]

\end{problemset}